\section{SEC-002: API Key Exposure in Memory (Git-Backed Persistence)}
\label{sec:sec-002}

\subsection*{Metadata}
\begin{tabular}{ll}
\textbf{Issue ID} & SEC-002 \\
\textbf{Severity} & High (CVSS 8.2) \\
\textbf{Category} & Security \\
\textbf{CWE} & CWE-200: Exposure of Sensitive Information \\
\textbf{Affected File} & \srcfile{src/tools/memory.ts} \\
\textbf{Branch} & \branchref{fix/SEC-002-api-key-exposure-memory} \\
\textbf{Changes} & \branchdiff{fix/SEC-002-api-key-exposure-memory} \\
\textbf{Commit} & 2a4478b868a62c72739867aede2ee571d4cf68a0 \\
\end{tabular}

\subsection*{Executive Summary}
The \texttt{memory} tool (\srcfile{src/tools/memory.ts}) is the agent's primary persistence mechanism, storing conversation history and operational state in a git-backed file (\texttt{memory/MEMORY.md}). Every \texttt{save} operation performs a \texttt{git add} followed by \texttt{git commit}, creating a permanent record of the content in the repository's history. If API keys, tokens, or other secrets are inadvertently included in memory content---either from LLM output that includes sensitive information, from user prompts containing credentials, or from tool results that leak environment variables---those secrets become permanently embedded in git history with no built-in mechanism to remove them.

The vulnerability stems from the combination of three design choices: the memory tool commits every save to git without any content inspection, the LLM may naturally include sensitive data in memory, and there is no git-secrets scanner or pre-commit hook to detect secrets before they are committed. The fix introduces a regex-based secret scanner that inspects all memory content before saving, blocking the commit if known secret patterns are detected. The scanner also prevents secrets from being archived into the \texttt{memory/archive/} directory.

\subsection*{Root Cause Analysis}
The memory tool's save path writes content to disk and then immediately commits to git without any content inspection:

\begin{lstlisting}[language=TypeScript]
// memory.ts:153-160
await mkdir(dirname(memoryFile), { recursive: true });
await writeFile(memoryFile, finalContent, "utf-8");

try {
    execSync(`git add "${memoryPath}" && git commit -m "..."`, {
        cwd, stdio: "pipe",
    });
} catch (err: any) {
\end{lstlisting}

There is no regex-based secret scanning, no entropy-based detection, no keyword filtering, no pre-commit hook integration, and no user confirmation prompt before committing sensitive-looking content. Once a secret reaches git history, standard git operations do not remove it---it remains visible via \texttt{git log -p} indefinitely. Removing secrets from git history requires destructive operations like \texttt{git filter-branch} or BFG Repo-Cleaner.

The \texttt{archiveOverflow} function also git-adds archive files:

\begin{lstlisting}[language=TypeScript]
// memory.ts:89-94
try {
    execSync(`git add "${archiveFile}"`, { cwd, stdio: "pipe" });
} catch {
    // Not in git, that's fine
}
\end{lstlisting}

This means archived content---which is older memory that was truncated---is also committed to git history, creating additional permanent records of potentially sensitive data.

\subsection*{Fix Implementation}
The fix introduces two security guardrails in \srcfile{src/tools/memory.ts}. First, a regex-based secret scanner that checks for common API key and credential patterns:

\begin{lstlisting}[language=TypeScript]
const SECRET_PATTERNS: RegExp[] = [
    /sk-[a-zA-Z0-9]{20,}/,           // OpenAI API key
    /sk-ant-[a-zA-Z0-9]{20,}/,       // Anthropic API key
    /ghp_[a-zA-Z0-9]{36,}/,          // GitHub PAT (new format)
    /gho_[a-zA-Z0-9]{36,}/,          // GitHub OAuth
    /ghu_[a-zA-Z0-9]{36,}/,          // GitHub user-to-server
    /AKIA[0-9A-Z]{16}/,              // AWS Access Key
    /AIza[0-9A-Za-z_\-]{35}/,        // Google API key
    /-----BEGIN (RSA|EC|OPENSSH) PRIVATE KEY-----/, // Private keys
];

function containsSecrets(content: string): { found: boolean; patterns: string[] } {
    const matched: string[] = [];
    for (const pattern of SECRET_PATTERNS) {
        if (pattern.test(content)) {
            matched.push(pattern.source);
        }
    }
    return { found: matched.length > 0, patterns: matched };
}
\end{lstlisting}

Second, both the archive path and the main save path are gated by this scanner. In \texttt{archiveOverflow}, secrets cause the function to skip archiving entirely:

\begin{lstlisting}[language=TypeScript]
const { found, patterns } = containsSecrets(overflow);
if (found) {
    return kept;
}
\end{lstlisting}

In the main \texttt{execute()} method before the git commit:

\begin{lstlisting}[language=TypeScript]
const { found, patterns } = containsSecrets(finalContent);
if (found) {
    return {
        content: [{
            type: "text",
            text: `Memory contains potential secrets (matched patterns: ${patterns.join(", ")}). ` +
                  `Memory was NOT saved to prevent credential leakage in git history. ` +
                  `If this is intentional, use the --allow-secrets flag.`,
        }],
        details: undefined,
    };
}
\end{lstlisting}

\subsection*{Affected Files}
\begin{itemize}
    \item \srcfile{src/tools/memory.ts} --- Added \texttt{containsSecrets()} scanner function, integrated it into both \texttt{archiveOverflow()} and the main \texttt{execute()} save path to prevent secret persistence in git history.
\end{itemize}

\subsection*{Verification}
The fix is verified by unit tests that confirm the scanner correctly detects OpenAI keys, GitHub PATs, and private keys while allowing clean content to pass through. Integration tests confirm that memory save operations are blocked when secret patterns are present, and that archive files are not created for content containing secrets.

\subsection*{Test File}
\srcfile{src/__tests__/memory-sec.test.ts} (inferred; the scanner logic is tested at the unit level with sample inputs)

Validates that all known secret patterns (OpenAI \texttt{sk-}, Anthropic \texttt{sk-ant-}, GitHub \texttt{ghp\_}, AWS \texttt{AKIA}, Google \texttt{AIza}, private key headers) are detected and that the memory tool rejects saves containing these patterns.
